\section*{الفصل \ref{ch:g7:central} --- التماثل المركزي ومتوازيات الأضلاع}

\begin{solution}{exo:g7:central:1}
صورة نصف الدورة: $M'$ على بُعد أربعة مربعات \emph{يسار} $O$ ومربع
واحد \emph{أسفل} (فالاتجاهان معكوسان معًا). وصورة الانعكاس
بالنسبة إلى الخط العمودي: أربعة مربعات يسارًا لكن ما زالت مربعًا واحدًا
\emph{أعلى}. والصورتان مختلفتان: فهما صورتان مرآتيتان عموديتان
إحداهما للأخرى.
\end{solution}

\begin{solution}{exo:g7:central:2}
$A'B' = AB$ (فأنصاف الدورة تحفظ الأطوال) و
$(A'B') \parallel (AB)$ (فصورة \omterm{def:g6:lines:objects}{المستقيم} \omterm{def:g6:lines:objects}{مستقيم} \omterm{def:g6:lines:perp}{موازٍ}،
\cref{prop:g7:central:preserved}).
\end{solution}

\begin{solution}{exo:g7:central:3}
أرقام آلة الحاسبة التي لها \omterm{def:g7:central:center}{مركز تماثل}: $0$ و $2$ و $5$ و $8$ (و
$1$ إذا رُسم عمودًا بسيطًا). والحروف: H و N و S و Z و O لها واحد؛ وأما A و
T فلا (لهما محور بدل ذلك).
\end{solution}

\begin{solution}{exo:g7:central:4}
الأضلاع المتقابلة متساوية: $MN = KL = 7$ سم و $NK = LM = 4$ سم. والزوايا
المتقابلة متساوية: $\widehat M = \widehat K = 115^\circ$. والقطران
يقطع كل منهما الآخر في منتصفهما المشترك، وتلك النقطة هي $O$: إذن
منتصف $[LN]$ هو $O$.
\end{solution}

\begin{solution}{exo:g7:central:5}
اُرسم \omterm{def:g6:lines:objects}{قطعة} $[AC]$ طولها $8$ سم، وعلّم منتصفها $O$، واُرسم مارًّا
بالنقطة $O$ مستقيمًا يصنع $60^\circ$ مع $(AC)$، وضع عليه $B$ و $D$
على بُعد $2.5$ سم من $O$ في الجهتين. وصِل $A$ و $B$ و $C$ و $D$: فالقطران
يتقاطعان في منتصفيهما، إذن $ABCD$ \omterm{def:g7:central:parallelogram}{متوازي أضلاع}
بالتعريف.
\end{solution}

\begin{solution}{exo:g7:central:6}
الزاويتان المتتاليتان متكاملتان: فبجانب $115^\circ$ تقع
$180 - 115 = 65^\circ$. والزاويتان المتقابلتان متساويتان، إذن الزوايا الأربع
هي $115^\circ$ و $65^\circ$ و $115^\circ$ و $65^\circ$ (\omterm{def:g1:addition:def}{والمجموع}
$360^\circ$).
\end{solution}

\begin{solution}{exo:g7:central:7}
من $A(1,1)$ إلى $O(2.5, 2.5)$: $1.5$ يمينًا و $1.5$ أعلى؛ وبمواصلة
الشيء نفسه: $A'(4, 4)$. ومن $B(4,2)$ إلى $O$: $1.5$ يسارًا و $0.5$ أعلى؛
وبالمواصلة: $B'(1, 3)$.
\end{solution}

\begin{solution}{exo:g7:central:8}
لا. فشبه المنحرف \omterm{def:g6:shapes:triangles}{المتساوي الساقين} ضلعاه المائلان متساويان دون
أن يكون \omterm{def:g7:central:parallelogram}{متوازي أضلاع} (فالضلعان المتوازيان لهما طولان
مختلفان). والذي يكفي هو المحكّ 3 من
\cref{thm:g7:central:recognize}: ضلعان متقابلان متساويان
\emph{ومتوازيان}.
\end{solution}

\begin{solution}{exo:g7:central:9}
النقطة $O$ منتصف $[BC]$ بالاختيار، وهي منتصف
$[AA']$ بإنشاء \omterm{def:g7:central:def}{المماثلة}. \omterm{def:g4:geometry:polygon}{والرباعي} $ABA'C$
قطراه $[AA']$ و $[BC]$ يشتركان في المنتصف $O$: فالمحكّ
1 (وهو التعريف) يجعله \omterm{def:g7:central:parallelogram}{متوازي أضلاع} بلا كلفة إضافية.
\end{solution}

\begin{solution}{exo:g7:central:10}
«الزوايا المتقابلة متساوية $\Rightarrow$ \omterm{def:g7:central:parallelogram}{متوازي أضلاع}»: \emph{صواب}
(وهو محكّ تعرّف آخر؛ فمع \omterm{def:g1:addition:def}{مجموع} الزوايا
$360^\circ$، يفرض تساوي الزوجين المتقابلين تكامل الزاويتين
المتتاليتين، ومنه توازي الأضلاع حسب
\cref{thm:g7:triangles:equalities}).

«القطران المتعامدان $\Rightarrow$ \omterm{def:g6:shapes:quadrilaterals}{معيّن}»: \emph{صواب} من أجل
\omterm{def:g7:central:parallelogram}{متوازي أضلاع} (فالقطران يتقاطعان أصلًا في منتصفيهما).

«\omterm{def:g3:shapes:rightangle}{زاوية قائمة} واحدة $\Rightarrow$ مستطيل»: \emph{صواب} ---
فالزاويتان المتتاليتان متكاملتان، إذن تصير الأربع قائمة.
\end{solution}

\begin{solution}{exo:g7:central:11}
مهما كان \omterm{def:g4:geometry:polygon}{الرباعي} --- حتى لو كان غير منتظم جدًّا --- تكوّن
المنتصفات $I$ و $J$ و $K$ و $L$ دائمًا \emph{\omterm{def:g7:central:parallelogram}{متوازي أضلاع}}.
والبرهان بمبرهنة المنتصفين في \cref{ch:g8:midpoints}
(\cref{exo:g8:midpoints:9}).
\end{solution}

\begin{solution}{pb:g7:central:1}
\textbf{1.} $M(3, 1) \to (-3, -1)$؛ و $N(-2, 4) \to (2, -4)$؛
و $P(0, -5) \to (0, 5)$. والقاعدة: نصف الدورة حول المبدأ
يرسل $(x, y)$ إلى $(-x, -y)$ --- فتتغيّر إشارة الإحداثيين معًا.

\textbf{2.} من $M(5, 3)$: ثلاثة مربعات يمين $C$ تصير
ثلاثة مربعات يسارًا، ومربعان أعلى يصيران مربعين أسفل: فالصورة
$(-1, -1)$. والتحقّق من الوصفة: ضعف المركز ناقص النقطة،
$(2 \times 2 - 5,\ 2 \times 1 - 3) = (-1, -1)$.

\textbf{3.} نصف الدورة حول مركز الورقة يبادل النقط
مثنى مثنى. وبعدد فردي من النقط، تبقى نقطة بلا شريك: فيجب
أن تكون صورة نفسها، والنقطة الوحيدة التي هي صورة نفسها
هي المركز نفسه. ومنه تقع النقطة الوسطى في الورقة 3 و5 و
7\,\dots\ في المركز تمامًا (اُنظر ورقة حقيقية: إنها كذلك).

\textbf{4.} حول $O_1(2,0)$: $A(1,1) \to (3,-1)$؛ وحول
$O_2(5,0)$: $(3,-1) \to (7,1)$. البداية $A(1,1)$، والنهاية
$(7,1)$: أي اِنزلاق $6$ إلى اليمين، لا قلب. والشيء نفسه من أجل
$B(0,3) \to (4,-3) \to (6,3)$: اِنزلاق $6$ يمينًا. فنصفا
الدورة يبلغان \emph{اِنسحابًا}.

\textbf{5.} الانزلاق $6$، و $O_1 O_2 = 3$: فالانسحاب
يقطع \emph{ضعف} المسافة بين المركزين (في
اِتجاه من $O_1$ إلى $O_2$) --- تمامًا كما اِنسحبت المرآتان المتوازيتان
بضعف الفجوة بينهما في \cref{pb:g6:symmetry:1}.

\textbf{6.} نصف الدورة حول $O$ يرسل \omterm{def:g7:central:parallelogram}{متوازي الأضلاع} على
نفسه (وهو ما يعنيه «\omterm{def:g7:central:center}{مركز التماثل}» له،
\cref{thm:g7:central:properties}). \omterm{def:g6:lines:objects}{والمستقيم} المارّ بالنقطة $O$ يُرسَل
أيضًا على نفسه، فتتبادل النقطتان اللتان يقطع فيهما
الحدّ موضعيهما: فهما متماثلتان بالنسبة إلى $O$. وتتبادل
قطعتا \omterm{def:g7:central:parallelogram}{متوازي الأضلاع} في جهتَي \omterm{def:g6:lines:objects}{المستقيم}
أيضًا، فهما نسختان مطابقتان (\cref{prop:g7:central:preserved})
--- ومتساويتا \omterm{def:g6:measure:area}{المساحة}.

\textbf{7.} اِقطع على \omterm{def:g6:lines:objects}{المستقيم} المارّ بالمركزين معًا:
مركز مستطيل الكعكة ومركز الثقب.
وحسب السؤال 6 مطبَّقًا على مستطيل الكعكة، ينصّف \omterm{def:g6:lines:objects}{القطع}
الكعكة؛ ومطبَّقًا على مستطيل الثقب (\omterm{def:g6:lines:objects}{فالمستقيم} يمرّ
بمركزه أيضًا)، ينصّف الثقب. فيأخذ كل جزء نصف الكعكة
ناقص نصف الثقب: أي عدل تامّ.

\textbf{8.} $D = (0 + 8 - 6,\ 0 + 5 - 1) = (2, 4)$ (فالقطران
يشتركان في منتصفهما، إذن $D$ صورة نصف الدورة للنقطة
$B$ حول المركز). والمركز: منتصف $[AC]$، أي
$O(4,\ 2.5)$.

\textbf{9.} في \omterm{def:g7:central:parallelogram}{متوازي الأضلاع} $ABA'C$، الأضلاع المتقابلة
متساوية: أي $BA' = AC$. والمتراجحة المثلثية في $ABA'$ تعطي
$AA' < AB + BA' = AB + AC$. وبما أن $O$ منتصف
$[AA']$، فإن $AA' = 2\,AO$، إذن $2\,AO < AB + AC$، أي\
$AO < \frac{AB + AC}{2}$.

\textbf{10.} الحدّ الأعلى: $AO < \frac{5 + 7}{2} = 6$ سم. والحدّ
الأدنى: في $ABA'$، $AA' > BA' - AB = 7 - 5 = 2$، إذن
$AO > 1$ سم. فالمتوسط الصادر من $A$ يقع تمامًا بين $1$ و
$6$ سم.

\textbf{11.} \omterm{def:g6:lines:objects}{القطعة}: نعم، منتصفها. \omterm{def:g6:lines:objects}{والمستقيم}: نعم --- فكل
نقطة من نقطه مركز. \omterm{def:g6:shapes:triangles}{والمثلث المتساوي الأضلاع}: لا مركز. والمربع:
نعم، تقاطع القطرين. والدائرة: نعم، مركزها.
والحرفان S و Z: نعم، نقطتهما الوسطى (اِقلب الصفحة رأسًا على عقب:
يُقرآن كما هما).

\textbf{12.} كان نصف الدورة يزاوج بين \omterm{def:g5:solids:def}{الرؤوس} الثلاثة
فيما بينها؛ وبما أن الثلاثة \omterm{def:g3:numbers:evenodd}{عدد فردي}، يكون رأس $A$ صورة
نفسه، فيُفرض أن يكون $A$ مركز نصف الدورة. ويكون الرأسان
الآخران $B$ و $C$ حينئذٍ متماثلين بالنسبة إلى $A$ ---
فيكون $A$ منتصف $[BC]$، فتقع \omterm{def:g5:solids:def}{الرؤوس} الثلاثة
على \omterm{def:g6:lines:objects}{مستقيم} واحد. وليس للمثلث ثلاثة \omterm{def:g5:solids:def}{رؤوس} مستقيمية:
فتناقض. إذن لا وجود لمركز.

\textbf{13.} حجّة الزوجية نفسها: فنصف دورة يحفظ
\omterm{def:g4:geometry:polygon}{المضلع} يزاوج بين \omterm{def:g4:geometry:polygon}{رؤوسه}؛ \omterm{def:g3:numbers:evenodd}{والعدد الفردي} من \omterm{def:g5:solids:def}{الرؤوس}
يفرض رأسًا ثابتًا، ولا بدّ أن يكون المركز ومنتصف
\omterm{def:g6:lines:objects}{القطعة} الواصلة بين رأسين آخرين --- أي ثلاثة \omterm{def:g5:solids:def}{رؤوس}
مستقيمية، وهذا مستحيل في \omterm{def:g4:geometry:polygon}{مضلع}\,\dots\ إذن لا \omterm{def:g4:geometry:polygon}{مضلع} ذو
\omterm{def:g3:numbers:evenodd}{عدد فردي} من \omterm{def:g5:solids:def}{الرؤوس} له \omterm{def:g7:central:center}{مركز تماثل}.

\textbf{14.} إنجاز نصف الدورة حول $O_1$، ثم حول
$O_2$، يرسل الشكل على نفسه في المرتين --- إذن
يفعل تركيبهما ذلك أيضًا. وحسب السؤال 5، ذلك التركيب
اِنسحاب بضعف المسافة $O_1 O_2$، وهو ليس معدومًا
لأن المركزين مختلفان. والشكل الذي يسع في ورقة لا يمكن أن
يساوي نفسه منزلقًا بمقدار ثابت غير معدوم (اِزلقه مرات
كافية فيغادر الورقة كلها). إذن للشكل المحدود
\omterm{def:g7:central:center}{مركز تماثل} واحد على الأكثر.

\textbf{15.} على الشريط اللانهائي، نصف دورة حول النقطة
الواقعة في منتصف ما بين أثر قدم يسرى وأثر القدم اليمنى التالية
يرسل النمط على نفسه؛ والنقطة الأبعد بخطوة كاملة
تفعل ذلك أيضًا: أي مركزان متمايزان. وتركيب نصفَي الدورة
يعطي الانسحاب بضعف نصف الخطوة --- وهو بالضبط
اِنزلاق الخطوة الواحدة الذي يرسل الإفريز اللانهائي بوضوح على
نفسه، كما يتوقّع السؤال 14. فالأنماط غير المحدودة تعيش
بقواعد مختلفة: وتلك هي رياضيات ورق الجدران.
\end{solution}
